\documentclass{eic-bcc-tcc}

%-----------------------------------------------
% Metadados do Trabalho
%-----------------------------------------------
\title{RageFilter: Classificação e Moderação Automática de Mensagens Tóxicas em Jogos \textit{ on-line }}
\titleen{RageFilter: Automatic Classification and Moderation of Toxic Messages in  on-line  Games}


% Nome completo do/da(s) autor(es)
% \addstudent[Full Name]{Nome Completo}
% O comando recebe um parâmetro opcional para a formatação dos nomes em inglÊs
\addstudent[Luiz Henrique Mendes de Oliveira Castro]{Luiz Henrique Mendes de Oliveira Castro}

\advisor{Marcelo Arêas Rodrigues da Silva, D.Sc.} 

% Membros da Banca para folha de aprovação
\addboardmember{Membro Interno}{Nome Membro Interno, D.Sc.\footnotemark[1]}
\addboardmember{Membro Externo}{Nome Membro Externo, D.Sc.\footnotemark[1] \footnotetext[1]{Caso a titulação não seja D.Sc., utilizar a adequada (p.ex., Dr. Ing., M.Sc., Esp., entre outras)}}

% Ano da defesa
\date{2025}


%-----------------------------------------------
% Abreviações
%-----------------------------------------------
% \newacronym{label}{sigla}{descrição} - Define um novo acrônimo
\newacronym{cefet}{Cefet/RJ}{Centro Federal de Educação Tecnológica Celso Suckow da Fonseca}
\newacronym{tcc}{TCC}{Trabalho de Conlusão de Curso}
\newacronym{abnt}{ABNT}{Associação Brasileira de Normas Técnicas}
\newacronym{ia}{IA}{Inteligência Artificial}


\begin{document}

%-----------------------------------------------
% Elementos pré-textuais
%-----------------------------------------------
\maketitle

\begin{dedicatoria}
Dedico este TCC aos meus pais e à minha namorada, por todo o amor, apoio e força que me acompanharam nesta caminhada.
\end{dedicatoria}


\begin{agradecimentos}
Agradeço, primeiramente, ao meu professor, cuja orientação, paciência e dedicação foram fundamentais para o desenvolvimento deste Trabalho de Conclusão de Curso. Sua capacidade de ensinar e incentivar o pensamento crítico foram essenciais para que eu pudesse evoluir ao longo desse processo.

Aos meus pais, deixo meu mais profundo agradecimento pelo amor incondicional, apoio constante e pela base sólida que sempre me proporcionaram. Sem eles, esta conquista não seria possível.

À minha namorada, agradeço pela compreensão nos momentos de ausência, pelo incentivo diário e por acreditar em mim mesmo mesmo quando eu duvidava. Sua parceria tornou esta jornada muito mais leve.

Ao meu primo, minha tia e minha avó, sou grato pelo apoio, carinho e pela torcida constante, sempre me lembrando da importância desta etapa.

Aos meus amigos da escola e da faculdade, agradeço pela parceria, pelas conversas, pela troca de conhecimentos e por tornarem o caminho mais divertido e menos solitário. Cada um de vocês contribuiu, de alguma forma, para que eu chegasse até aqui.

A todos que, direta ou indiretamente, fizeram parte desta trajetória, deixo minha sincera gratidão.

Por fim, agradeço ao CEFET/RJ, instituição que me acompanha desde 2021 e que foi fundamental na minha formação acadêmica, oferecendo estrutura, conhecimento e oportunidades que contribuíram profundamente para o meu crescimento pessoal e profissional.
\end{agradecimentos}

\begin{resumo}
O crescimento das comunidades digitais de jogos eletrônicos tem gerado interações tóxicas em plataformas como Twitch, Reddit e Discord. Apesar de estudos sobre o tema, ainda faltam \textit{datasets} em português focados na linguagem \textit{gamer}, o que dificulta a criação de ferramentas de moderação automática adequadas. Este trabalho propõe o desenvolvimento de um sistema de detecção e censura de mensagens tóxicas, utilizando técnicas de Processamento de Linguagem Natural. A metodologia envolve a coleta e anotação de mensagens, criação de um dataset especializado e treinamento de modelos transformer para identificar e moderar a toxicidade. O sistema também terá um mecanismo de censura seletiva. A solução visa promover um ambiente digital mais seguro para jogadores, com contribuições parciais já alcançadas, como a definição da taxonomia de toxicidade e o planejamento do \textit{pipeline} computacional.

\palavraschave{Toxicidade em jogos. Processamento de Linguagem Natural. Modelo de classificação. Moderação automática. Comunidades gamer.}
\end{resumo}


\begin{abstract}
\textit{The growth of digital gaming communities has led to toxic interactions on platforms like Twitch, Reddit, and Discord. Despite studies on the topic, there is still a lack of Portuguese datasets focused on gamer language, which hinders the development of adequate automatic moderation tools. This work proposes the development of a system for detecting and censoring toxic messages using Natural Language Processing techniques. The methodology involves collecting and annotating messages, creating a specialized dataset, and training transformer models to identify and moderate toxicity. The system will also feature a selective censorship mechanism. The solution aims to create a safer digital environment for gamers, with partial contributions already achieved, such as defining the toxicity taxonomy and planning the computational pipeline.}

\keywords{Toxicity detection. Natural Language Processing. Classification model. Automatic moderation. Gaming communities.}
\end{abstract}

% Lista de acrônimos
%\listofacronymspage{LIS}

% Sumário
\tableofcontentspage

% --- ELEMENTOS TEXTUAIS ---
\mainmatter

\chapter{Introdução}
\label{cha:introducao}

O crescimento da indústria dos jogos digitais tem transformado plataformas como Twitch\footnote{https://www.twitch.tv/}, Reddit\footnote{https://www.reddit.com/}, Discord\footnote{https://discord.com/} e fóruns especializados em espaços centrais de socialização, competição e construção de comunidades on-line. No entanto, junto com esse avanço, observou-se também a intensificação de comportamentos tóxicos, linguagem agressiva e práticas de assédio dirigidas a jogadores, \textit{streamers} e participantes dessas comunidades, como podemos observar no estudo feito em \cite{gam3sgg}. A presença constante desse tipo de mensagem compromete a experiência dos usuários e levanta discussões sobre bem-estar digital, moderação de conteúdo e criação de ambientes mais seguros dentro do ecossistema dos \textit{games}.

Ao mesmo tempo, modelos de inteligência artificial vêm se consolidando como ferramentas eficazes para análise textual, detecção de padrões de linguagem e filtragem de conteúdos potencialmente ofensivos. A combinação entre a necessidade de moderação automatizada e o avanço de técnicas de aprendizado de máquina abre espaço para pesquisas que desenvolvam soluções específicas para o universo gamer, um contexto marcado por gírias próprias, expressões rápidas, abreviações e xingamentos pouco estudados em bases tradicionais de toxicidade. Nesse cenário, torna-se relevante investigar como mensagens tóxicas se manifestam dentro das comunidades de jogos e como modelos podem ser treinados para reconhecê-las com precisão.

Diante desse panorama, este trabalho busca contribuir para o desenvolvimento de tecnologias de moderação automática capazes de atuar em ambientes digitais relacionados a \textit{games}, construindo um \textit{dataset} especializado e treinando modelos voltados à detecção de toxicidade textual. Além de preencher lacunas existentes na literatura, especialmente no contexto brasileiro, a pesquisa oferece subsídios para iniciativas que desejam mitigar comportamentos hostis, melhorar a convivência on-line e promover espaços mais saudáveis para interação entre jogadores.

\section{Motivação}

O ambiente dos jogos on-line é caracterizado por interações rápidas e espontâneas, frequentemente marcadas por emoções intensas. Comunidades e grupos de jogos competitivos se tornaram grandes centros de comunicação digital, mas também locais onde insultos, ataques pessoais e comportamentos tóxicos se proliferam com facilidade. Esse fenômeno tem impacto direto na experiência dos usuários, podendo afastar jogadores, influenciar negativamente o bem-estar psicológico e comprometer a permanência em comunidades virtuais.

Embora diversas plataformas adotem ferramentas de moderação, grande parte desses sistemas utiliza modelos generalistas que não capturam nuances linguísticas específicas do vocabulário \textit{gamer}, como gírias, abreviações e xingamentos próprios desse contexto. Além disso, há escassez de \textit{datasets} públicos dedicados exclusivamente ao estudo da toxicidade em jogos, o que limita o desenvolvimento de modelos especializados. Essa lacuna motiva o desenvolvimento do presente trabalho.

\section{Definição do Problema}

O desafio central deste TCC consiste em identificar e classificar automaticamente mensagens potencialmente tóxicas presentes em comunidades on-line relacionadas a jogos digitais, considerando suas particularidades linguísticas e contextuais. Para isso, é necessário coletar, organizar e rotular mensagens oriundas de plataformas como Twitch, Reddit e grupos de jogos, além de compilar um conjunto específico de palavras e expressões de baixo calão frequentemente utilizadas no cenário \textit{gamer}.

Dessa forma, o problema pode ser enunciado da seguinte maneira: como desenvolver um modelo de aprendizado de máquina capaz de identificar toxicidade textual em mensagens de comunidades \textit{gamer}, utilizando um \textit{dataset} especializado composto por mensagens reais e por um vocabulário próprio desse ambiente?

A resposta a esse problema exige tanto a construção do \textit{dataset} quanto o desenvolvimento e avaliação do modelo de classificação, compondo a essência deste trabalho.

\section{Objetivos}

Geral: Desenvolver um modelo de aprendizado de máquina capaz de identificar e censurar mensagens tóxicas em comunidades on-line relacionadas a jogos digitais, utilizando um \textit{dataset} próprio construído a partir de mensagens reais e vocabulário específico do ambiente \textit{gamer}.

Específicos:
\begin{itemize}
    \item Coletar mensagens provenientes de plataformas como Twitch, Reddit e outras comunidades  on-line  ligadas ao universo \textit{gamer}.
    \item Construir um \textit{dataset} de toxicidade contendo mensagens rotuladas como tóxicas e não tóxicas.
    \item Elaborar um segundo \textit{dataset} contendo palavras, gírias e xingamentos típicos do ambiente \textit{gamer}, especialmente em língua portuguesa.
    \item Realizar a limpeza, padronização e pré-processamento dos textos coletados.
    \item Avaliar diferentes modelos de aprendizado supervisionado para classificação de toxicidade.
    \item Treinar o modelo selecionado utilizando os \textit{datasets} construídos.
    \item Analisar o desempenho obtido por meio de métricas adequadas, como precisão, recall e F1-score.
    \item Propor formas de aplicação prática do modelo em sistemas de moderação automática.
\end{itemize}

\section{Organização dos Capítulos}

O presente trabalho está dividido da seguinte forma: Capítulo 2, apresenta a Fundamentação Teórica. O Capítulo 3, Desenvolvimento, expondo a proposta e a solução. Por fim, o Capítulo 4, Avaliação Experimental, que aponta os resultados atingidos e os próximos passos para o projeto.

\chapter{Fundamentação Teórica}
\label{cha:fundamentacao}

Este capítulo apresenta os conceitos fundamentais para compreender o fenômeno da toxicidade em comunidades de jogos digitais, o funcionamento de chats multiplayer, bem como os métodos computacionais utilizados para detecção automática de mensagens ofensivas.

\section{Jogos Multiplayers com Chat}

Jogos \textit{multiplayer}  on-line  consolidaram-se como espaços de interação social intensa, nos quais jogadores cooperam, competem e dialogam em tempo real. Diferentemente de mídias tradicionais, esses ambientes combinam estímulos emocionais, anonimato e dinâmicas de competição, fatores que moldam o comportamento dos participantes.

Segundo \cite{sbgames}, plataformas como Twitch e Discord funcionam como extensões sociais do ambiente de jogo, possibilitando conversas síncronas e assíncronas, acompanhando transmissões ao vivo, discutindo estratégias e reforçando o sentimento de comunidade. Os autores destacam que a Twitch concentra cerca de 67\% do mercado global de \textit{streaming} ao vivo e movimenta mais de 2,3 milhões de usuários ativos, cuja comunicação textual ocorre em ritmo acelerado e com linguagem própria do meio digital. 

O Discord, por sua vez, organiza comunidades em servidores temáticos, possibilitando canais permanentes de bate-papo por voz e texto. Essa infraestrutura torna as interações contínuas, mesmo fora do contexto imediato da partida, aumentando a produção textual e, consequentemente, a possibilidade de ocorrência de conflitos, insultos e hostilidade.

Essas características fazem dos chats \textit{multiplayer}, tanto internos ao jogo quanto externos, um objeto de estudo relevante para análise de comportamentos linguísticos e para o desenvolvimento de modelos automáticos de moderação.

\section{Toxicidade em Comunidades de Jogos  on-line }

A toxicidade em jogos digitais refere-se ao conjunto de comportamentos verbais e atitudinais que prejudicam a experiência de outros jogadores. Isso inclui insultos, xingamentos, assédio, discriminação racial ou de gênero, ameaças e incitação ao ódio.

Para compreender e visualizar a toxicidade nos jogos, \cite{assali2024comportamento} apresentam uma análise detalhada da toxicidade em League of Legends\footnote{https://www.leagueoflegends.com/pt-br/}, mostrando que insultos relacionados à habilidade, como \textit{"noob"}, "inútil", "lixo", os mais recorrentes, seguidos de xingamentos agressivos e comportamentos discriminatórios. O estudo categoriza diferentes tipos de insultos e mostra que mais de 90\% dos jogadores já testemunharam comportamentos tóxicos dentro do jogo, enquanto cerca de 40\% admitem ter sido responsáveis por tais atitudes em algum momento .

Segundo \cite{sbgames}, ataques direcionados a mulheres, minorias raciais e pessoas LGBTQIA+ são comuns em chats de Twitch e Discord, reforçados pela cultura competitiva, anonimato e pela normalização de agressões verbais nas comunidades gamer. Esses comportamentos podem provocar impactos psicológicos negativos, como estresse, insegurança e exclusão social .

Em síntese, a toxicidade é um fenômeno complexo, multifatorial e linguístico, demandando abordagens computacionais robustas para sua identificação.

\section{Modelos de Linguagem para Análise e Classificação de Toxicidade}

O crescimento no volume de dados textuais em jogos e plataformas de \textit{streaming} torna inviável a moderação totalmente humana. Para lidar com isso, têm sido utilizados modelos de Processamento de Linguagem Natural (PLN) e Large Language Models (LLMs) na detecção de toxicidade.

Modelos como BERT, RoBERTa, T5 e variantes multilíngues são amplamente empregados para classificação textual, incluindo detecção de toxicidade, discurso de ódio e assédio. Eles se destacam pela capacidade de aprender representações profundas da linguagem, lidar com ironias, gírias, abreviações e compreender dependências contextuais.

No trabalho \cite{naseem-etal-2025-gametox}, os autores demonstram que abordagens conjuntas de \textit{intent classification} (classificação da intenção da frase) e \textit{slot filling} (anotação \textit{token} a \textit{token}) produzem melhores resultados do que modelos que utilizam apenas rótulos em nível de sentença. O \textit{dataset} criado pelos autores contém 53 mil mensagens de \textit{World of Tanks}\footnote{https://worldoftanks.com/pt-br/}, rotuladas em categorias como \textit{Insults and Flaming} (Insultos e 'Farpas'), \textit{Hate and Harassment} (Ódio e Assédio), \textit{Threats} (Ameaças) e \textit{Extremism} (Extremismo), além de \textit{slots} específicos como \textit{Game Slang} (Gírias de Jogos) e \textit{Toxic}.

Essa abordagem aumenta a granularidade da análise e melhora a generalização, pois o modelo aprende tanto o contexto da frase quanto as palavras específicas responsáveis pela toxicidade.

Essas evidências reforçam o uso de modelos baseados em \textit{embeddings} contextuais e arquiteturas transformer para o desenvolvimento de soluções automatizadas de censura e classificação.

\section{Trabalhos Relacionados}

\subsection{Metodologia da busca}

A metodologia de busca adotada para a identificação dos trabalhos relacionados teve como objetivo mapear estudos relevantes que abordassem a toxicidade em jogos digitais e em plataformas adjacentes, bem como pesquisas que utilizassem técnicas computacionais para análise automática de linguagem ofensiva. A estratégia de busca foi conduzida de forma exploratória e direcionada, buscando equilibrar abrangência temática e pertinência ao escopo deste trabalho.

Inicialmente, foram definidos critérios de inclusão baseados nos principais eixos do TCC. Foram considerados trabalhos que abordassem ao menos um dos seguintes aspectos: análise de toxicidade em jogos on-line; estudos sobre comportamento tóxico em plataformas relacionadas aos jogos, como serviços de \textit{streaming} e fóruns; criação ou utilização de \textit{datasets} voltados à detecção de toxicidade; aplicação de técnicas de Processamento de Linguagem Natural para identificação de discurso ofensivo; e uso de modelos de linguagem em contextos associados a comunidades \textit{gamer}. Trabalhos que não apresentavam relação direta com jogos digitais ou que tratavam de toxicidade em contextos muito distintos, como ambientes corporativos ou políticos, foram excluídos.

A busca bibliográfica foi realizada nas bases de dados Scopus\footnote{https://www.periodicos.capes.gov.br/?option=com_pcollection&mn=70&smn=79&cid=63}
, IEEE Xplore\footnote{https://ieeexplore.ieee.org/Xplore/home.jsp}
 e Google Scholar\footnote{https://scholar.google.com.br}
, além de canais de conferências e periódicos da área. As buscas utilizaram operadores booleanos \textit{AND} e \textit{OR}, combinando termos relacionados à toxicidade, jogos digitais e métodos computacionais. Termos equivalentes foram agrupados com \textit{OR}, enquanto conceitos distintos foram combinados com \textit{AND}.

Foram utilizadas múltiplas \textit{strings} de busca, tanto em inglês quanto em português, uma vez que cada idioma retorna conjuntos distintos de resultados. Exemplos em inglês incluem: \textit{("game toxicity" OR "toxic behavior in online games") AND ("natural language processing" OR "toxicity detection")}. Em português, utilizaram-se combinações como: \textit{("toxicidade em jogos online" OR "comportamento tóxico em jogos") AND ("processamento de linguagem natural" OR "modelos de linguagem")}.

Após a recuperação inicial dos trabalhos, foi realizada uma triagem baseada na leitura dos títulos, resumos e palavras-chave, com o objetivo de identificar estudos alinhados ao escopo do projeto. Em seguida, os artigos selecionados passaram por uma leitura mais aprofundada, permitindo avaliar sua contribuição, metodologia e resultados. Trabalhos duplicados, excessivamente generalistas ou sem fundamentação metodológica clara foram descartados.

Por fim, os estudos selecionados foram organizados de acordo com sua abordagem principal, distinguindo-se entre pesquisas de caráter sociocultural e estudos com foco computacional. Essa organização permitiu comparar diferentes perspectivas sobre a toxicidade em comunidades \textit{gamer} e identificar lacunas na literatura, especialmente no que diz respeito à ausência de \textit{datasets} em português brasileiro e de soluções automáticas voltadas à moderação de conteúdo.

\subsection{Apresentação dos trabalhos relacionados}

a) \cite{assali2024comportamento}

Este estudo qualitativo, conduzido por Pires e Fortim (2021), realiza uma análise detalhada de toxicidade no jogo League of Legends, investigando a frequência de insultos e comportamentos agressivos, bem como os impactos psicológicos sobre os jogadores. Os autores categorizaram diferentes tipos de ofensas, como "insultos diretos", "xenofobia" e "misoginia", e examinaram como fatores como o anonimato nas plataformas de jogos contribuem para a intensificação dessas atitudes. O estudo também discutiu as implicações sociais desses comportamentos, como o estresse, a exclusão e a insegurança gerados nas vítimas de ataques verbais, além de explorar as dinâmicas sociais de grupos de jogadores.

Contribuição: O trabalho fornece uma compreensão sociocultural profunda sobre os fatores que motivam a toxicidade dentro de jogos on-line, especialmente League of Legends, ao mostrar que o comportamento tóxico não é apenas um problema de comportamento individual, mas um reflexo de dinâmicas sociais mais amplas, como a normalização da violência no contexto digital e o papel do anonimato.

b) \cite{sbgames}

O artigo de Nunes e Fortim (2025) amplia a compreensão sobre a toxicidade nas plataformas adjacentes aos jogos on-line, como Twitch e Discord. Este estudo investiga interações tóxicas em \textit{chats} ao vivo e em discussões em fóruns, destacando como o ambiente de \textit{streaming} e as características sociais dessas plataformas contribuem para o aumento dos discursos discriminatórios, como racismo, misoginia e homofobia. Os autores analisam o papel da comunicação síncrona (ao vivo) e assíncrona, destacando que a falta de moderação eficaz nessas plataformas, aliada ao anonimato e à alta participação, favorece a propagação desses comportamentos.

Contribuição: O estudo é significativo pois amplia o entendimento de toxicidade para além do contexto dos jogos em si, considerando o papel das plataformas de \textit{streaming} e fóruns como arenas adicionais onde a toxicidade se manifesta. Ele também destaca os efeitos do anonimato e da normalização de discursos agressivos no contexto dessas plataformas digitais.


c) \cite{naseem-etal-2025-gametox}

O estudo de \cite{naseem-etal-2025-gametox} apresenta um \textit{dataset} de 53 mil mensagens coletadas do jogo \textit{World of Tanks}, rotuladas com diferentes categorias de toxicidade, incluindo \textit{Insults} and \textit{Flaming}, \textit{Hate} and \textit{Harassment}, \textit{Threats} e \textit{Extremism}, além de slots específicos como \textit{Game Slang} e \textit{Toxic}. Os autores exploram diferentes modelos de PLN, incluindo variantes de BERT, para classificar as mensagens, e demonstram que abordagens conjuntas de \textit{intent classification} (classificação da intenção da mensagem) e \textit{slot filling} (anotação \textit{token} a \textit{token}) fornecem melhores resultados em termos de acurácia e generalização.

Contribuição: Este trabalho oferece uma base valiosa de dados para a criação de modelos de PLN, além de apresentar uma abordagem mista de classificação de intenções e análise detalhada do conteúdo das mensagens. A pesquisa é útil para o desenvolvimento de modelos que consideram não apenas o texto completo, mas também as palavras-chave responsáveis pela toxicidade.


d) \cite{doi:10.1177/13548565211028809}

Gandolfi e Ferdig  apresentam uma análise sobre a toxicidade nas plataformas de serviço de jogos (\textit{Game Service Platforms} - GSPs), com foco específico no jogo DOTA 2\footnote{https://www.dota2.com/home} e nas plataformas Twitch.tv e Steam. O estudo foi realizado com dados coletados ao longo de quatro semanas, envolvendo chats ao vivo, \textit{reviews} de usuários e discussões em fóruns. Os autores utilizam uma abordagem de análise de discurso para compreender como a toxicidade se manifesta nessas plataformas e como ela é legitimada, reforçada ou desafiada pelos participantes. A pesquisa discute os diferentes tipos de toxicidade observados, incluindo agressões raciais, homofóbicas e sexistas, e explora como a hegemonia masculina e a competitividade exacerbada contribuem para a manutenção de comportamentos tóxicos. Além disso, o estudo destaca como as características das plataformas, como a comunicação síncrona em Twitch e a interatividade mais assíncrona em Steam, influenciam a dinâmica e a propagação da toxicidade.

Contribuição: A pesquisa amplia o entendimento de como a toxicidade se manifesta em plataformas além do jogo em si, trazendo \textit{insights} sobre o comportamento dos jogadores e os fatores sociais que facilitam a toxicidade.

\chapter{Desenvolvimento}
\label{cha:desenvolvimento}

Este capítulo descreve a proposta de solução, os artefatos que compõem o sistema, bem como o processo de construção dos \textit{datasets} e dos modelos computacionais destinados à detecção de toxicidade em mensagens oriundas de plataformas relacionadas ao universo gamer. O objetivo é apresentar a modelagem conceitual da solução, os componentes planejados, o fluxo geral do sistema e as etapas de desenvolvimento que serão consolidadas no TCC II.

\section{Descrição da Proposta}

A proposta deste trabalho consiste no desenvolvimento de um sistema computacional capaz de identificar e censurar mensagens tóxicas em ambientes digitais associados ao universo gamer, especialmente plataformas como Twitch, Reddit, Discord e chats internos de jogos. A ideia central é criar uma solução baseada em técnicas de Processamento de Linguagem Natural (PLN) que consiga lidar com a linguagem altamente informal, dinâmica e repleta de gírias característica desses espaços. Para isso, o sistema partirá da construção de um \textit{dataset} original em português brasileiro, composto por mensagens coletadas nessas plataformas e rotuladas de acordo com diferentes categorias de toxicidade identificadas na literatura e adaptadas ao contexto nacional. A partir desse \textit{dataset}, será treinado um modelo de classificação textual capaz de analisar novas mensagens, detectar conteúdos ofensivos e aplicar mecanismos de censura seletiva ou integral, dependendo da gravidade da toxicidade detectada.

A proposta se justifica pela ausência de bases de dados abrangentes em português brasileiro focadas especificamente na linguagem gamer, o que dificulta o uso de modelos generalistas e reduz a eficiência de sistemas automáticos de moderação. Ao criar um \textit{dataset} especializado e um modelo treinado com esses dados, o trabalho pretende contribuir tanto para a área acadêmica quanto para aplicações práticas de moderação automática em comunidades virtuais de jogos.

\section{Modelagem da Solução}

A modelagem da solução envolve a definição dos módulos conceituais que compõem o sistema e das relações entre eles. O primeiro módulo é o da coleta de dados, que será responsável por reunir mensagens reais provenientes das plataformas selecionadas. Essa etapa exigirá métodos de extração distintos para cada fonte, já que a Twitch oferece APIs e arquivos exportáveis de transmissões, o Reddit possui ferramentas como o Pushshift\footnote{https://pushshift.io/} para recuperar comentários e Praw\footnote{https://praw.readthedocs.io/en/stable/} para ler postagens, enquanto o Discord depende da autorização de servidores para exportação dos registros. Toda essa coleta deve seguir princípios éticos e procedimentos de anonimização para garantir a privacidade dos usuários envolvidos.


Concluída a coleta, inicia-se a etapa de preparação e construção do \textit{dataset}. Todo o conteúdo será consolidado em um único repositório, passando por padronização de codificação, limpeza de ruídos, remoção de duplicatas e formatação uniforme. A seguir, será realizada a rotulação das mensagens, em um processo dividido em duas camadas complementares. Na primeira, cada mensagem será classificada de acordo com sua natureza e intensidade de toxicidade. Na segunda, será feita uma anotação com \textit{tokens}, na qual palavras específicas serão identificadas como tóxicas, gírias \textit{gamer}, verbos ou elementos neutros, seguindo inspiração de \textit{datasets} como o GameTox\footnote{https://github.com/shucoll/GameTox}. Essa estrutura detalhada permitirá que o modelo aprenda tanto o contexto geral da mensagem quanto os elementos linguísticos que caracterizam a toxicidade, fortalecendo sua capacidade de generalização.


Além do \textit{dataset} principal, será desenvolvido um segundo repositório contendo exclusivamente gírias, expressões, abreviações e insultos característicos do cenário \textit{gamer} brasileiro. Esse conjunto funcionará como um suporte para o processo de anotação e poderá ser incorporado ao modelo como conhecimento adicional, possibilitando que ele reconheça termos que não costumam aparecer em bases de dados generalistas.

O modelo de classificação será treinado utilizando arquiteturas do tipo transformer, como BERTimbau\footnote{https://github.com/neuralmind-ai/portuguese-bert/}, mBERT\footnote{https://github.com/google-research/bert} ou variantes do LLaMA\footnote{https://www.llama.com} adaptadas para português. Esses modelos possuem capacidade de capturar relações contextuais profundas, o que é fundamental para identificar nuances da toxicidade, especialmente em mensagens curtas e informais. O treinamento envolverá etapas com \textit{tokens}, normalização, divisão dos dados em conjuntos de treino e teste, ajuste de hiperparâmetros e avaliação com métricas apropriadas como precisão, recall e F1-score. Como a toxicidade costuma aparecer de forma desbalanceada — com maior volume de mensagens não tóxicas — será necessário aplicar técnicas de balanceamento para evitar que o modelo seja enviesado.

Por fim, será elaborado o mecanismo de censura, responsável por aplicar a moderação automática às mensagens analisadas pelo modelo. Esse componente atuará traduzindo a classificação em uma ação prática. Para mensagens extremamente ofensivas, a censura poderá ocultar o texto por completo. Para mensagens com toxicidade moderada ou localizada, a censura seletiva permitirá substituir apenas as palavras ofensivas por símbolos, preservando o restante do conteúdo. Esse tipo de abordagem só é possível graças à anotação com \textit{tokens} descrita anteriormente, que identifica precisamente os trechos problemáticos da mensagem.

\begin{figure}[htb]
    \centering
    \includegraphics[width=0.8\textwidth]{Fluxograma.png}
    \caption{Diagrama das etapas da Modelagem da Solução}
    \label{fig:arquitetura}
\end{figure}


A modelagem apresentada estabelece a arquitetura geral da solução e evidencia a forma como os módulos se integram ao longo do desenvolvimento do trabalho. Conforme sintetizado na Figura~\ref{fig:arquitetura}, as etapas representam o encadeamento lógico das atividades planejadas e executadas até o momento. Cada decisão técnica foi concebida de modo a permitir que o sistema opere de forma eficiente sobre a linguagem \textit{gamer} brasileira, preservando simultaneamente princípios éticos, consistência metodológica e viabilidade computacional.

\section{Construção da Solução}

Como este trabalho corresponde ao TCC I, a construção da solução encontra-se em estágio inicial, concentrando-se nas definições conceituais e no delineamento metodológico. Até o momento, foram realizados estudos aprofundados sobre toxicidade em jogos digitais, compreensão das características linguísticas das plataformas escolhidas, levantamento de requisitos e desenho da arquitetura geral do sistema. A partir desses estudos, tornou-se possível definir a taxonomia inicial de categorias de toxicidade e o método de anotação a ser adotado no \textit{dataset}. 

Atualmente, como principais fontes de coleta de dados, foram escolhidas três \textit{lives} na plataforma Twitch de dois grandes \textit{streamers} do jogo League of Legends (Jukes e Bahiano). As \textit{lives} foram escolhidas de modo aleatório e delas, foi coletado apenas o \textit{chat} ao vivo das gravações. Totalizando aproximadamente quarenta mil frases enviadas por usuários. Porém, a grande maioria dos textos, era composto por \textit{emotes}, risadas, funcionalidades da Twitch e até \textit{spam} de letras ou siglas, então foi realizada uma limpeza, reduzindo o número para aproximadamente quinze mil.

Além disso, foram extraídos da rede social reddit, dos canais brasileiros relacionados aos jogos League of Legends, Counter-Strike\footnote{https://www.counter-strike.net/cs2} 2, DOTA 2 e Overwatch\footnote{https://overwatch.blizzard.com/pt-br/} as quinhentas publicações em alta e as quinhentas publicações mais recentes, totalizando em torno de onze mil retornos. 

No TCC II, serão executadas as etapas práticas de implementação. A coleta de dados será concluída com a consolidação das mensagens provenientes de Twitch, Reddit, Discord e outras fontes relevantes. O \textit{dataset} será construído e anotado integralmente, seguindo tanto a classificação por frase quanto a marcação com \textit{tokens}. Em seguida, os modelos de linguagem serão treinados e avaliados conforme o plano estabelecido. A implementação do mecanismo de censura integrará os módulos anteriores, permitindo que o sistema funcione de forma completa e autônoma. Por fim, todo o processo experimental será documentado, fornecendo métricas e análises que demonstrarão a eficácia da solução desenvolvida.

O desenvolvimento apresentado neste capítulo constitui, portanto, a base conceitual da solução. Ele estabelece as etapas, os módulos e os princípios necessários para que, no TCC II, a implementação possa ocorrer de forma estruturada, coerente e alinhada aos objetivos gerais da pesquisa.

\chapter{Avaliação Experimental}
\label{cha:avaliacao}

\section{Planejamento da Validação}

Como este trabalho se encontra em sua fase inicial, a validação experimental será realizada integralmente no TCC II. Ainda assim, é possível delinear, nesta etapa, o planejamento conceitual da avaliação, descrevendo os métodos previstos, os critérios de análise e o procedimento experimental a ser seguido. O objetivo da validação será mensurar, de forma objetiva, o desempenho do modelo de detecção de toxicidade a partir de métricas amplamente adotadas em Processamento de Linguagem Natural, além de avaliar qualitativamente a adequação do sistema ao contexto \textit{gamer} brasileiro.

A validação deverá ocorrer em duas dimensões complementares. A primeira dimensão será a avaliação quantitativa do modelo, buscando medir sua capacidade de identificar corretamente mensagens tóxicas e não tóxicas, bem como distinguir entre diferentes categorias de toxicidade. A segunda será uma avaliação qualitativa do uso prático do mecanismo de censura seletiva, verificando se a substituição de palavras ou a ocultação completa de mensagens ocorre de maneira coerente com o contexto de uso real.

\section{Procedimento Experimental Previsto}

O experimento previsto iniciará pela divisão do \textit{dataset} em três subconjuntos: treino, validação e teste. Essa divisão permitirá avaliar o modelo com dados inéditos, garantindo que os resultados reflitam sua capacidade real de generalização. O modelo será treinado com o conjunto de treino e terá seus hiperparâmetros ajustados com o conjunto de validação. Em seguida, será aplicado ao conjunto de teste, no qual será possível mensurar seu desempenho por meio de métricas como precisão, recall, F1-score e acurácia. Como a toxicidade tende a apresentar desbalanceamento entre classes, métricas macro-ponderadas e análises específicas por categoria também serão utilizadas para evitar que o desempenho aparente seja influenciado pelo predomínio de mensagens não tóxicas.

Na etapa qualitativa, a avaliação envolverá a observação direta das saídas do sistema, especialmente do módulo de censura. Serão analisados casos em que apenas parte da mensagem é censurada e casos em que a toxicidade detectada justifica a censura completa. Essa análise permitirá identificar possíveis falhas do modelo, como censuras excessivas, falta de censura em trechos problemáticos ou interpretações incorretas do contexto. A intenção é compreender, além das métricas numéricas, como o sistema se comporta diante de mensagens reais do ambiente gamer, preservando elementos essenciais da comunicação sem comprometer a segurança do usuário.

\subsection{Critérios de Avaliação}

Os critérios adotados para a avaliação serão definidos de forma a refletir tanto o rigor científico quanto as necessidades práticas de um sistema de moderação. Na dimensão quantitativa, a escolha das métricas seguirá padrões consolidados na literatura, garantindo comparabilidade com outros estudos e permitindo análises detalhadas por tipo de toxicidade. Na dimensão qualitativa, a avaliação se concentrará na coerência das decisões do sistema, na naturalidade da censura realizada e em sua adequação à linguagem real das plataformas analisadas. Dessa forma, espera-se obter uma visão mais ampla e precisa do desempenho da solução.


\chapter{Considerações Parciais}

O desenvolvimento realizado até o presente momento permitiu estruturar a base conceitual e metodológica necessária para a construção do sistema de detecção de toxicidade proposto. Os objetivos específicos definidos na Introdução orientaram cada etapa conduzida no TCC I, e todos os esforços realizados até aqui convergem para permitir, no TCC II, a execução prática e experimental do projeto.

O primeiro objetivo específico consistia em compreender o funcionamento das plataformas e dos ambientes digitais de onde os dados serão coletados. Esse objetivo foi atingido por meio da revisão de literatura e da análise detalhada de plataformas como Twitch, Reddit e Discord, que revelou suas características linguísticas, sociais e estruturais. O segundo objetivo dizia respeito à definição da taxonomia de toxicidade e do modelo teórico de rotulação, o que também foi alcançado com base nos estudos realizados e na adaptação de categorias presentes em trabalhos como GameTox, ajustando-as ao contexto da língua portuguesa. A partir desses estudos, tornou-se possível construir o desenho preliminar do \textit{dataset}, definindo tanto o nível de anotação por sentença quanto a anotação por \textit{token} que será empregada no TCC II.

O terceiro objetivo específico envolvia a modelagem do pipeline computacional para detecção de toxicidade. Essa etapa foi concluída com a definição dos módulos da solução, incluindo coleta de dados, construção de \textit{dataset}, treinamento de modelos de PLN e implementação do mecanismo de censura. Embora não tenha sido realizada nenhuma implementação prática até o momento, a arquitetura conceitual encontra-se completamente delineada. Por fim, o quarto objetivo, que trata da preparação para o processo de validação experimental, foi parcialmente atendido, uma vez que o planejamento da avaliação encontra-se descrito, ainda que sua execução prática seja reservada para o TCC II.

Até o momento, não houve resultados experimentais ou modelos treinados, pois esta etapa pertence integralmente ao desenvolvimento futuro. Entretanto, o trabalho já possui uma base sólida, com diretrizes claras sobre quais dados deverão ser coletados, como serão anotados, quais modelos poderão ser empregados e como será conduzida a fase de avaliação. O progresso realizado no TCC I garante que o TCC II poderá se concentrar na implementação efetiva da solução, na execução dos experimentos e na análise quantitativa e qualitativa dos resultados.

Para que o projeto seja concluído integralmente, ainda falta realizar a coleta completa das mensagens, construir e anotar o \textit{dataset}, treinar e ajustar os modelos de classificação textual, implementar o mecanismo de censura com base nos rótulos de toxicidade e conduzir a validação experimental definida no capítulo anterior. Esses elementos constituem a espinha dorsal do TCC II e representam o núcleo prático do trabalho.

\section{Cronograma}

O cronograma do Trabalho de Conclusão de Curso foi elaborado com o objetivo de organizar as etapas do desenvolvimento da pesquisa ao longo do período de agosto a dezembro.

Nos meses de agosto e setembro, foram realizadas pesquisas preliminares para definição e aprofundamento do tema, bem como o levantamento das tecnologias a serem utilizadas. Nesse período, também se iniciou a redação das seções iniciais do trabalho.

Em outubro, as atividades concentraram-se na pesquisa de fontes e referências bibliográficas e no desenvolvimento da proposta inicial, contribuindo para a avaliação da viabilidade do projeto.

No início do mês de novembro, após o avanço das pesquisas e do desenvolvimento, constatou-se que a proposta inicial, que previa o desenvolvimento de um jogo, não atendia plenamente aos objetivos do TCC. Diante disso, optou-se pela mudança do tema, buscando maior alinhamento com o escopo do trabalho.

A partir dessa redefinição, as atividades de redação, desenvolvimento e revisão foram intensificadas ao longo dos meses de novembro e dezembro, culminando na revisão final do trabalho e na preparação para a apresentação à banca examinadora.

\begin{table}[h!]
\centering
\small
\setlength{\tabcolsep}{4pt}
\begin{tabular}{lcccccc}
\hline
\textbf{Atividade} & \textbf{Ago} & \textbf{Set} & \textbf{Out} & \textbf{Nov} & \textbf{Dez} & \textbf{Jan} \\ \hline
Pesquisas para o TCC              & X & X &   & X & X &   \\
Mudança de Tema de TCC            &   &   &   & X &   &   \\
Pesquisa de Tecnologias Adequadas & X & X &   &   &   &   \\
Redação                           & X & X &   & X & X &   \\
Pesquisa de Fontes e Referências  &   &   & X & X &   &   \\
Desenvolvimento                   &   &   & X & X &   &   \\
Revisão e Ajustes                 &   &   &   &   & X &   \\
Apresentação à Banca              &   &   &   &   &   & X \\ \hline
\end{tabular}
\caption{Cronograma de atividades realizadas ao longo da construção do TCC 1}
\end{table}


\section{Próximos Passos}

As próximas etapas, que serão realizadas ao longo do TCC II, envolvem a construção concreta dos artefatos projetados. A primeira delas é a coleta definitiva dos dados em todas as plataformas planejadas, seguida da organização, limpeza e padronização desses dados. Em seguida, será feita a anotação das mensagens, utilizando tanto rotulação manual quanto estratégias de apoio automático, possibilitando a criação de um \textit{dataset} robusto e representativo. Após essa fase, será iniciado o treinamento dos modelos de PLN, com rodadas sucessivas de ajuste de parâmetros e comparações de desempenho. Na etapa final, será implementado o mecanismo de censura, integrando-o ao modelo treinado, e então será executada a validação experimental que deverá medir a eficácia do sistema. Todo esse processo será documentado rigorosamente para garantir reprodutibilidade e transparência científica.

\begin{table}[h!]
\centering
\resizebox{\textwidth}{!}{%
\begin{tabular}{lcccccc}
\hline
\textbf{Atividade} & \textbf{Fev} & \textbf{Mar} & \textbf{Abr} & \textbf{Mai} & \textbf{Jun} & \textbf{Jul} \\ \hline
Redação                               & X & X & X & X & X & X \\
Desenvolvimento                       & X & X & X & X & X &   \\
Pesquisa para aplicações              &   & X & X & X &   &   \\
Levantamento de testes e aplicações   &   &   & X & X &   &   \\
Revisão e Ajustes                     &   &   &   & X & X & X \\
Apresentação à Banca                  &   &   &   &   &   & X \\ \hline
\end{tabular}
}
\caption{Cronograma de Planejamento para desenvolvimento do TCC2}
\end{table}

\section{Uso de Inteligência Artificial}

Durante o desenvolvimento deste trabalho, foi utilizada a ferramenta de inteligência artificial ChatGPT, da OpenAI, para auxiliar na elaboração da redação e no desenvolvimento de parte do código. O ChatGPT foi utilizado como suporte para a redação de seções do texto, aprimoramento de ideias, formatação de conteúdo e construção de tabelas, além de auxiliar no desenvolvimento e depuração de código, especialmente nas seções relacionadas à implementação técnica e algoritmos.

% --- ELEMENTOS PÓS-TEXTUAIS ---
\backmatter

% % Referências
\referencepage


\end{document}